% References (BGU \S 10.7.2.10), Section 8.
% Hand-formatted IEEE-style reference list. The entries are typed in IEEE
% format directly (rather than generated by a CSL/.bst) so the same numbered,
% IEEE-styled list renders identically in the XeLaTeX proof PDF and in the
% Pandoc/DOCX deliverable. \bibitem keys match every \cite key in the body;
% numbering follows order of first citation.
\section{References}
\label{sec:references}
\renewcommand{\refname}{}

\begin{thebibliography}{99}

\bibitem{he2016resnet}
K.~He, X.~Zhang, S.~Ren, and J.~Sun, ``Deep residual learning for image
recognition,'' in \emph{Proc. IEEE Conf. Comput. Vis. Pattern Recognit.
(CVPR)}, 2016, pp. 770--778.

\bibitem{huang2017densenet}
G.~Huang, Z.~Liu, L.~van der Maaten, and K.~Q. Weinberger, ``Densely connected
convolutional networks,'' in \emph{Proc. IEEE Conf. Comput. Vis. Pattern
Recognit. (CVPR)}, 2017, pp. 2261--2269.

\bibitem{shi2024transnext}
D.~Shi, ``TransNeXt: Robust foveal visual perception for vision transformers,''
arXiv:2311.17132, 2024.

\bibitem{huynh2008psnr}
Q.~Huynh-Thu and M.~Ghanbari, ``Scope of validity of PSNR in image/video
quality assessment,'' \emph{Electron. Lett.}, vol. 44, no. 13, pp. 800--801,
2008.

\bibitem{wang2004ssim}
Z.~Wang, A.~C. Bovik, H.~R. Sheikh, and E.~P. Simoncelli, ``Image quality
assessment: From error visibility to structural similarity,'' \emph{IEEE Trans.
Image Process.}, vol. 13, no. 4, pp. 600--612, 2004.

\bibitem{hendrycks2021naturaladv}
D.~Hendrycks, K.~Zhao, S.~Basart, J.~Steinhardt, and D.~Song, ``Natural
adversarial examples,'' in \emph{Proc. IEEE/CVF Conf. Comput. Vis. Pattern
Recognit. (CVPR)}, 2021.

\bibitem{krizhevsky2009cifar10}
A.~Krizhevsky, ``Learning multiple layers of features from tiny images,'' Univ.
of Toronto, Tech. Rep., 2009.

\bibitem{lecun1998mnist}
Y.~LeCun, L.~Bottou, Y.~Bengio, and P.~Haffner, ``Gradient-based learning
applied to document recognition,'' \emph{Proc. IEEE}, vol. 86, no. 11, pp.
2278--2324, 1998.

\bibitem{loshchilov2019adamw}
I.~Loshchilov and F.~Hutter, ``Decoupled weight decay regularization,'' in
\emph{Proc. Int. Conf. Learn. Represent. (ICLR)}, 2019.

\bibitem{loshchilov2017sgdr}
I.~Loshchilov and F.~Hutter, ``SGDR: Stochastic gradient descent with warm
restarts,'' in \emph{Proc. Int. Conf. Learn. Represent. (ICLR)}, 2017.

\bibitem{naeini2015ece}
M.~P. Naeini, G.~F. Cooper, and M.~Hauskrecht, ``Obtaining well calibrated
probabilities using Bayesian binning,'' in \emph{Proc. AAAI Conf. Artif.
Intell.}, 2015.

\bibitem{guo2017oncalibration}
C.~Guo, G.~Pleiss, Y.~Sun, and K.~Q. Weinberger, ``On calibration of modern
neural networks,'' in \emph{Proc. Int. Conf. Mach. Learn. (ICML)}, 2017.

\bibitem{huang2016stochastic}
G.~Huang, Y.~Sun, Z.~Liu, D.~Sedra, and K.~Q. Weinberger, ``Deep networks with
stochastic depth,'' in \emph{Proc. Eur. Conf. Comput. Vis. (ECCV)}, 2016.

\bibitem{zhang2018mixup}
H.~Zhang, M.~Cisse, Y.~N. Dauphin, and D.~Lopez-Paz, ``mixup: Beyond empirical
risk minimization,'' in \emph{Proc. Int. Conf. Learn. Represent. (ICLR)}, 2018.

\bibitem{yun2019cutmix}
S.~Yun, D.~Han, S.~J. Oh, S.~Chun, J.~Choe, and Y.~Yoo, ``CutMix:
Regularization strategy to train strong classifiers with localizable
features,'' in \emph{Proc. IEEE/CVF Int. Conf. Comput. Vis. (ICCV)}, 2019.

\bibitem{falcon2019lightning}
W.~Falcon \emph{et al.}, ``PyTorch Lightning,'' GitHub repository,
Lightning-AI/lightning, 2019.

\bibitem{paszke2019pytorch}
A.~Paszke \emph{et al.}, ``PyTorch: An imperative style, high-performance deep
learning library,'' in \emph{Adv. Neural Inf. Process. Syst. (NeurIPS)}, vol.
32, 2019.

\end{thebibliography}
